%\documentclass{Math_Note}
%\setlength{\parindent}{2em}

%\title{\textbf{Mathmatical Analysis Zorich Comprehensions Chapter 3}}
%\author{Buce-Ithon}
%\newdateformat{mydate}{\twodigit{\THEDAY}{ }\shortmonthname[\THEMONTH], \THEYEAR}
%\date{\today}

%\begin{document}
% Title page
%\maketitle

% Content
%\newpage
%\tableofcontents
%\newpage

% \setcounter{section}{-1}
\newpage
\section{Chapter 3.Limits}
\subsection{Comprehensions}
\subsubsection{$\xi.1$ The Limit of a Sequence}
First, the most basic is the $\varepsilon$-$N$ definition of sequence limits and the basic properties of limits (existence $\implies$ boundedness). 
On this basis, the operational processes of limits are established (four arithmetic operations and inequality relations, note: the strictness of 
inequalities may not be maintained after the limit process). Then, we proceed to the discussion of the existence of limits (introducing two types 
of convergence criteria: the /general/ Cauchy convergence criterion and the /special - monotonic sequence/ Weierstrass theorem), and an important 
method: subsequences, subsequence limits, and upper and lower limits.

Furthermore, this section also introduces some preliminary knowledge of series. The essential idea is to treat series as partial sum sequences, so 
that the theory of sequence limits above can be directly applied to obtain convergence tests for series (also divided into general and special two 
types) and introduces absolute convergence as a stronger form of convergence (unaffected by rearrangement), and the related properties of convergent 
series (necessary condition - terms $\to 0$). On this basis (mainly the Weierstrass theorem), some useful corollaries are made, such as: the comparison 
test (only applicable when all terms are non-negative or all terms have the same sign), the Cauchy convergence criterion (based on the relationship 
between $\sup \lim_{n\to\infty} |a_n|^{1/n}$ and $1$) and the d'Alembert test (based on the relationship between $\sup \lim_{n\to\infty} |\frac{a_{n+1}}{a_n}|$ 
and $1$) (applicable in general cases and ensures absolute convergence), and the Cauchy condensation test ($\sum_{n=1}^{\infty} a_n$ and 
$\sum_{k=0}^{\infty} 2^k a_{2^k}$ have the same convergence behavior) (applicable when terms are non-negative and decreasing).
\newline\newline
\fbox{\begin{minipage}{0.95\textwidth}
\textbf{Examples}
\begin{enumerate}
    \item[1.] For the series $\sum_{n=1}^{\infty} (2 + (-1)^n)^n x^n$, it converges for $|x| < \frac{1}{3}$ and diverges for $|x| > \frac{1}{3}$ and 
    $|x| = \frac{1}{3}$, using the Cauchy root test.
    \item[2.] For the series $\sum_{n=1}^{\infty} \frac{x^n}{n!}$, it converges for all $x \in \mathbb{R}$, using the d'Alembert ratio test. Note: This 
    can also be judged by the Cauchy root test, based on $\lim_{n \to +\infty} \sqrt[n]{n!} = +\infty$.
    \item[3.] For the series $\sum_{n=1}^{\infty} \frac{1}{n^p}$, it converges for $p > 1$ and diverges for $p \leq 1$, using the Cauchy condensation test.
\end{enumerate}
\end{minipage}}

\subsubsection{$\xi.2$ The Limit of a Function}
The study of function limits largely follows the methods of sequence limits. First, the $\varepsilon$-$\delta$ definition of a function limit is introduced 
(note that function limits are defined on a punctured neighborhood), along with the properties of a function's limit existing at a point (existence $\implies$ 
local boundedness). On this basis, the four arithmetic operations of function limits are studied (the proof process will become much more convenient if an 
infinitesimal quantity is first introduced as an equivalent definition of the limit of function) and the inequality relations during the limit process (note: 
the strictness of inequalities may not be maintained after the limit process). Furthermore, the limit behavior of elementary functions and the limit behavior of 
composite functions are explained (here, it is necessary to learn two important limits). 

A novel aspect of Zorich's approach is his definition of something called a "base" (essentially, a family of sets satisfying certain conditions, which facilitates 
the description of "approaching" in the limit process); he then re-describes the aforementioned content based on this concept. Next, the discussion proceeds to 
the problem of the existence of limits (again introducing two methods: the /general/ Cauchy convergence criterion and the /special - monotonic function/ Weierstrass 
theorem). Finally, the issue of asymptotic comparison of functions (order estimation) is briefly introduced.
\newline\newline
\fbox{\begin{minipage}{0.95\textwidth}
\textbf{2 Important Limits}
\begin{enumerate}
    \item[1.] $\lim_{x \to 0} \frac{\sin x}{x} = 1$.
    \item[2.] $\lim_{x \to \infty} \left(1+\frac{1}{x}\right)^x = e$.
\end{enumerate}
\end{minipage}}

\newpage
\subsection{Problems Sets}
\subsubsection{$\xi.1$ The Limit of a Sequence}
\begin{prb}
    (C3S1Q3)\\
    We mark all the points on a circle obtained from a fixed point by rotations of the circle through angles of $n$ radians, where $n \in \mathbb{Z}$ ranges over 
    all integers. Describe all the limit points of the set so constructed.
\end{prb}
\fbox{\begin{minipage}{0.95\textwidth}
\begin{sol*}
Attempt 1. $n(mod\ \alpha)=n-\alpha\left\lfloor  \frac{n}{\alpha}  \right\rfloor,\ \forall\alpha\in\mathbb{R}$. \\
Attempt 2. $\forall\varepsilon>0$, by pigeonhole principle ($\forall N>0, as\ n\to\infty$, $n-N$ of integers ($n-N\to\infty$) will be located between those $N$ integers.) 
$\exists m,n< \frac{1}{\varepsilon}$, s.t. $|m-n|(mod\ 2\pi)<\varepsilon$, \\
then let $k=|m-n|$, (construct full proof for dense property), $\forall x\in[0,2\pi]$, $x\in\left[ k\left\lfloor \frac{x}{k(mod\ 2\pi)} \right\rfloor,k\lceil \frac{x}{k(mod\ 2\pi)} \rceil \right]$ 
and the length of interval equals to $k<\varepsilon$, which means $\mathbb{R}_{2\pi}$ is dense in $(0,2\pi)$ or $[0,2\pi]$.
\end{sol*}
\end{minipage}}
\newline\newline
\fbox{\begin{minipage}{0.95\textwidth}
\textcolor{skycyan}{Development:} \\
Kronecker's approximation theorem.(Wolfram MathWorld, Wikipedia)
\end{minipage}}
\newline\newline
(Variants) \\
Prove: $\{n^p\pmod{\pi}:p>0\}$ is dense on $[0,\pi]$.
\newline\newline
\fbox{\begin{minipage}{0.95\textwidth}
\begin{pf*}
Core: Similar to the idea of differential (tool: difference), reduce the order, and finally turn it into the prototype of the above problem. \\
Let $p>0$ be an arbitrary real number. We claim that $\{n^p\pmod{\pi}:p>0\}$ is dense on $[0,\pi]$. \\
For any function $f(n)$ let $\Delta f(n)=f(n+1)-f(n)$. If $f(n)$ converges to $0\pmod{\pi}$ then so does $\Delta f(n)$. Now for $f(n)=n^p$ we have 
$\Delta f(n) \approx pn^{p-1}$, $\Delta^2 f(n) \approx p(p-1)n^{p-2}$, etc. More precisely, by Newton's binomial theorem, $(n+1)^p-n^p=\sum_{i=1}^{\infty} \tbinom{p}{i}n^{p-i}$. \\
Let $k = \lfloor p \rfloor$. One can check that the part of the series with $i > k$ has the overall form $o(1)$. Thus $\Delta (n^p) = pn^{p-1} + c_2 n^{p-2} + \cdots + c_k n^{p-k} + o(1)$ 
for various constant coefficients $c_2, \ldots, c_k$. Iterating, it follows that $\Delta^k (n^p) = (p)_k n^{p-k} + o(1)$, where $(p)_k = p(p - 1) \cdots (p - k + 1)$. Thus the function 
$g(n) = (p)_k n^{p-k}$ tends to $0$ mod $\pi$, but this is impossible: since $\pi$ is irrational and we must have $p > k$, but then $g(n) \to \infty$ while $\Delta g(n) \to 0$ (There aren't 
$\pmod{\pi}$, which means $g(n)$ tends to $\infty$ "densely".), so the range of $g$ is dense mod $\pi$.
\end{pf*}
\end{minipage}}

\begin{prb}
    (C3S1Q5)
    \begin{enumerate}
        \item[a)] $1+\frac{1}{1!}+\frac{1}{2!}+\cdots+\frac{1}{n!}+\frac{1}{n!n}=3-\frac{1}{1\cdot 2\cdot 2!}-\cdots-\frac{1}{(n-1)\cdot n\cdot n!}$ holds, as $n\geq 2$.
        \item[b)] $e=3-\sum_{n=0}^{\infty} \frac{1}{(n+1)\cdot (n+2)\cdot (n+2)!}$.
        \item[c)] For computing $e$,$e\approx 1+\frac{1}{1!}+\frac{1}{2!}+\cdots+\frac{1}{n!}+\frac{1}{n!n}$ is much better than $e\approx 1+\frac{1}{1!}+\frac{1}{2!}+\cdots+\frac{1}{n!}$.
    \end{enumerate}
\end{prb}
\fbox{\begin{minipage}{0.95\textwidth}
\begin{qsol*}
\begin{enumerate}
    \item[a)] $1 + \frac{1}{1!} + \frac{1}{2!} + \cdots + \frac{1}{n!} + \frac{1}{n!n} + \frac{1}{1\cdot 2\cdot 2!} + \cdots + \frac{1}{(n-1)\cdot n\cdot n!}$ 
    $=\left[ \frac{1}{n!} + \frac{1}{n!n} + \frac{1}{(n-1)n\cdot n!} - \frac{1}{(n-1)\cdot (n-1)!} \right]$ 
    $=1 + \frac{1}{1!} + \frac{1}{2!} + \cdots + \frac{1}{(n-1)!} + \frac{1}{(n-1)\cdot (n-1)!} + \frac{1}{1\cdot 2\cdot 2!} + \cdots + \frac{1}{(n-2)\cdot (n-1)\cdot (n-1)!}$ 
    $=\cdots = 1 + \frac{1}{1!} + \frac{1}{1\cdot 1!} = 3$
    \item[b)] $\lim_{n\to\infty} 1 + \frac{1}{1!} + \frac{1}{2!} + \cdots + \frac{1}{n!} + \frac{1}{n!n}$ and 
    $\lim_{n\to\infty} 1 + \frac{1}{1!} + \frac{1}{2!} + \cdots + \frac{1}{n!}$, $\Delta_n = \frac{1}{n!n} = o(1), n\to\infty.$
    \item[c)] Since both $1 + \frac{1}{1!} + \frac{1}{2!} + \cdots + \frac{1}{n!} + \frac{1}{n!n}$ and $1 + \frac{1}{1!} + \frac{1}{2!} + \cdots + \frac{1}{n!}$ are monotonically increasing 
    $\to e$, and the former has an extra term $\frac{1}{n!n}$ thus it grows faster (i.e., converges faster). To estimate the difference in convergence speed, we just need to find the maximum 
    $k \in \mathbb{N}^+$ such that $\frac{1}{n!n} \geq \frac{1}{(n+1)!} + \cdots + \frac{1}{(n+k)!}$.\\
    (Actually, $\forall n \geq 2, k \in \mathbb{N}^+$ i.e $Max{k} = +\infty$)
\end{enumerate}
\end{qsol*}
\end{minipage}}

\begin{prb}
    (C3S1Q6) \\
    Given$S_p(a, b) = \left( \frac{a^p + b^p}{2} \right)^{\frac{1}{p}}, p \neq 0, a, b > 0$, prove: 
    \begin{enumerate}
        \item[a)] $\forall p \neq 0, a, b > 0, S_p(a, b) \in [a, b]$;
        \item[b)] $\lim_{n \to \infty} S_n(a, b) = b, \lim_{n \to \infty} S_{-n}(a, b) = a$.
    \end{enumerate}
\end{prb}
\fbox{\begin{minipage}{0.95\textwidth}
\begin{qsol*}
\begin{enumerate}
    \item[a)] It is sufficient to define the order relation of $a$ and $b$, for example: without loss of generality, assume $a < b$.
    \item[b)] It is sufficient to define the relative size relationship of $a$ and $b$ (i.e., what is commonly called substitution or elimination by increment), for example: without loss of 
    generality, let $b = \lambda a$.
\end{enumerate}
\end{qsol*}
\end{minipage}}
\newline\newline
\fbox{\begin{minipage}{0.95\textwidth}
\textcolor{skycyan}{Thoughts:} \\
(\textcolor{cyan}{Trick}) When a problem has multiple ($\geq 2$) variables, it may be helpful to look for connections between these seemingly "independent" variables 
(for example, in question 6.b) of this section).
\end{minipage}}

\begin{prb}
    (C3S1Q7) (Classical problem) \\
    If $a>0$, then sequence $x_{n+1} = \frac{1}{2}\left(x_n + \frac{a}{x_n}\right)$ converges to $\sqrt{a}$.   
\end{prb}
\fbox{\begin{minipage}{0.95\textwidth}
\begin{qsol*}
$x_{n+1}=\frac{1}{2}\left( x_{n}+ \frac{a}{x_{n}} \right) \geq \sqrt{a},\ n \geq 1$, i.e. $\forall x_{1}>0,\ x_{2} \geq \sqrt{a}$. \\
$[$Induction$]$ then $\{x_{n}\}_{n=2}^{\infty}$ monotonically decrease, $\lim_{n\to\infty} x_{n+1}=\lim_{n\to\infty} \frac{1}{2}\left( x_{n}+\frac{1}{x_{a}} \right)=\sqrt{a}$.
\end{qsol*}
\end{minipage}}
\newline\newline
\fbox{\begin{minipage}{0.95\textwidth}
\textcolor{skycyan}{Thoughts:} \\
(\textcolor{cyan}{Experience}) An experienced student, upon seeing a recurrence relation, can immediately guess that the problem will most likely require the use of methods like the Weierstrass theorem to 
first determine convergence, and then take the limit on both sides of the equation. \\
\end{minipage}}
\newline\newline
\fbox{\begin{minipage}{0.95\textwidth}
(\textcolor{cyan}{Tip}) If starting from a general problem-solving approach, the first step is usually to determine the basic properties of the sequence (positivity/negativity, monotonicity, parity, periodicity 
(or these properties starting from a certain term), etc.), and then consider which method to use to address convergence (general → Cauchy convergence criterion, monotonic → Weierstrass theorem) 
and other issues.
\end{minipage}}
\newline\newline
(Follow-up Question) \\
Please estimate the convergence rate, i.e., the dependence of $x_n - \sqrt{a}$ on $n$.

\begin{prb}
    (C3S1Q8)
    \begin{enumerate}
        \item[a)] (Faulhaber's formula) $S_{k}(n)=\sum_{i=1}^{n} i^k=\frac{1}{k+1}\sum_{r=0}^{k}C_{k+1}^{r}B_{r}n^{k+1-r}$, $B_{r}$ is Bernoulli number.
        \item[b)] $\lim_{ n \to \infty } \frac{S_{k}(n)}{n^{k+1}} = \frac{1}{k+1}$.
    \end{enumerate}
\end{prb}
\fbox{\begin{minipage}{0.95\textwidth}
\begin{sol*}
\begin{enumerate}
    \item[a)] This part is a little hard, because to wholy solving it,  we needs expotential generating function with indeterminate $z$: $G(z,n)=\sum_{k=0}^{\infty}S_{k}(n) \frac{1}{k!} z^k$. You could get it from wikipedia.
    \item[b)] Method1. $\lim_{ n \to \infty } \frac{S_{k}(n)}{n^{k+1}} = \lim_{ n \to \infty } \frac{1}{n}\sum_{i=1}^{n} \frac{i^k}{n^k} = \int_{0}^{1} x^k dx = \frac{1}{k+1}$ (Riemann sum) \\
    Method2. $\lim_{ n \to \infty } \frac{S_{k}(n)}{n^{k+1}} = \lim_{ n \to \infty } \frac{(n+1)^k}{(n+1)^{k+1}-n^{k+1}}$ (Stolze-Cesaro Lemma) $=\lim_{ n \to \infty } \frac{1}{(n+1)-\left( \frac{n}{n+1} \right)^k\cdot n} 
    = \lim_{ n \to \infty } \frac{1}{1+n\left[ 1-\left( \frac{n}{n+1} \right)^k \right]}$ (order estimation) $=\lim_{ n \to \infty } \frac{1}{1+n\left[ 1-\left( 1-k\cdot \frac{1}{n+1}+O\left( \frac{1}{n^{2}} \right) \right) 
    \right]} = \lim_{ n \to \infty } \frac{1}{1+ \frac{n}{n+1}\cdot k+O\left( \frac{1}{n} \right)}=\frac{1}{1+k}$. 
\end{enumerate}
\end{sol*}
\end{minipage}}

\subsubsection{$\xi.2$ The Limit of a Function}
\begin{prb}
    (C3S2Q1) \\
    Porve that: 
    \begin{enumerate}
        \item[a)] there $\exists$ a unique function defined on $\mathbb{R}$ satisfying : 
        $f(1) = a$ (where $a > 0, a \neq 1$), $f(x_1) \cdot f(x_2) = f(x_1 + x_2)$, $\lim_{x \to x_0} f(x) = f(x_0)$.
        \item[b)] there $\exists$ a unique function defined on $\mathbb{R}^+$ satisfying : 
        $f(a) = 1$ (where $a > 0, a \neq 1$), $f(x_1) + f(x_2) = f(x_1 \cdot x_2)$, $\lim_{x \to x_0} f(x) = f(x_0)$.
    \end{enumerate}
\end{prb}
\fbox{\begin{minipage}{0.95\textwidth}
\begin{qsol*}
By substitution (for example: $x \triangleq e^{\xi}$, etc.), transform the given variant of the Cauchy equation back into the form $f(x + y) = f(x) + f(y)$ (uniqueness is obtained from the equation itself and continuity, and 
the method has been provided in the discussion of the corresponding equation in Fikhtengoltz's "Calculus Tutorial (I)").
\end{qsol*}
\end{minipage}}

\begin{prb}
    (C3S2Q2) \\
    \begin{enumerate}
        \item[a)] Provide an isomorphism mapping between the group $(\mathbb{R}, +)$ and the group $(\mathbb{R}^+, \cdot)$;
        \item[b)] Prove that $(\mathbb{R}, +)$ and the group $(\mathbb{R} \setminus {0}, \cdot)$ are not isomorphic.
    \end{enumerate}
\end{prb}
\fbox{\begin{minipage}{0.95\textwidth}
\begin{qsol*}
\begin{enumerate}
    \item[a)] $\forall x, y \in \mathbb{R}$, if $\varphi(x + y) = \varphi(x) \cdot \varphi(y)$, then $\varphi(x) = e^{cx}$;
    \item[b)] From the fact that "isomorphism not only is a bijection but also preserves the group operation", if $\varphi: \mathbb{R} \to \mathbb{R} \setminus {0}$ exists, then from $\varphi(0) = 1$, it follows that 
    $\varphi(x) \cdot \varphi(-x) = \varphi(0) = 1$.
\end{enumerate}
\end{qsol*}
\end{minipage}}

\begin{prb}
    (C3S2Q3)
\end{prb}
\fbox{\begin{minipage}{0.95\textwidth}
\begin{qsol*}
\begin{enumerate}
    \item[a)] $\lim\limits_{ x \to +0 }x^{x}=\lim\limits_{ x \to +0 }e^{ x\ln(x) }=\lim\limits_{ t \to +\infty }e^{ -\frac{\ln(t)}{t} }=1$; 
    \item[b)] $\lim\limits_{ x \to +\infty } x^{\frac{1}{x}} = \lim\limits_{ x \to +\infty }e^{\frac{\ln(x)}{x}} = 1$; 
    \item[c)] $\lim\limits_{ x \to 0 } \frac{\log_{a}(1+x)}{x} = \lim\limits_{ x \to 0 } \frac{1}{\ln a} \cdot \frac{\ln(1+x)}{x}$ $= \left[ \ln(1+x)=x-\frac{1}{2}x^{2}+O(x^{3}) \right] = \frac{1}{\ln a}$; 
    \item[d)] $\lim\limits_{ x \to 0 } \frac{a^x-1}{x} = [a^x-1 \triangleq t] = \lim\limits_{ x \to 0 } \frac{t}{\log_{a}(1+t)} = \ln a$. 
\end{enumerate}
\end{qsol*}
\begin{nt*}
Both a) and b) used $\lim_{ x \to +\infty } \frac{\log_{a}x}{x^{\alpha}} = 0,\ \alpha>0$.
\end{nt*}
\end{minipage}}

\begin{prb}
    (C3S2Q4) \\
    Prove: $1+\frac{1}{2}+\cdots+\frac{1}{n}=\ln n+c+o(1),\ n\to \infty$.
\end{prb}
\fbox{\begin{minipage}{0.95\textwidth}
\begin{qsol*}
Using $\ln(n+1)-\ln n=\ln \frac{n+1}{n}=\ln\left( 1+\frac{1}{n} \right)=\frac{1}{n}+O\left( \frac{1}{n^{2}} \right),\ n\to \infty$, then $\Sigma$ them.
\end{qsol*}
\end{minipage}}

\begin{prb}
    (C3S2Q5)
    \begin{enumerate}
        \item[a)] If two positive term series $\sum_{n=1}^{\infty} a_n$ and $\sum_{n=1}^{\infty} b_n$ satisfy: $a_n \sim b_n, n \to \infty$, then these two series have the same convergence behavior;
        \item[b)] The series $\sum_{n=1}^{\infty} \sin \frac{1}{n^p}$ converges only when $p > 1$.
    \end{enumerate}
\end{prb}
\fbox{\begin{minipage}{0.95\textwidth}
\begin{sol*}
\begin{enumerate}
    \item[a)] $a_n \sim b_n$ i.e. $\forall \varepsilon > 0, \exists N > 0, s.t. \forall n > N, |\frac{a_n}{b_n} - 1| < \varepsilon$ i.e. 
    $1 - \varepsilon < \frac{a_n}{b_n} < 1 + \varepsilon, \forall n > N, \text{take a specific } \varepsilon = \varepsilon_0, \text{then}$ 
    $(1 - \varepsilon_0)b_n < a_n < (1 + \varepsilon_0)b_n. \text{Convergence/divergence is thereby controlled.}$
    \item[b)] $p > 0 : \sin \frac{1}{n^p} \sim \frac{1}{n^p}$ Two positive term series have the same convergence behavior, the convergence/divergence of the latter can be given by the Cauchy condensation test 
    (see Summary C3S1 for details) \\ $p = 0 : \sin \frac{1}{n^p} = \sin 1 > 0$ diverges \\ $p < 0 : \sum_{n=1}^\infty \sin \frac{1}{n^p} = [q \triangleq -p > 0] = \sum_{n=1}^\infty \sin n^q$ diverges (Essentially, 
    $\{n^q (\text{mod } \pi)\}$ is dense on $[0, \pi]$, details in (C3S1Q3 variants).
\end{enumerate}
\end{sol*}
\end{minipage}}

\begin{prb}
    (C3S2Q6)
    \begin{enumerate}
        \item[a)] If $\forall n \in \mathbb{N}$ we have $a_n \geq a_{n+1} > 0$, and the series $\sum_{n=1}^{\infty} a_n$ converges, then $a_n = o\left(\frac{1}{n}\right)$, $n \to \infty$;
        \item[b)] If $b_n = o\left(\frac{1}{n}\right)$, then it is always possible to construct a convergent series $\sum_{n=1}^{\infty} a_n$; such that $b_n = o(a_n)$, $n \to \infty$;
        \item[c)] If the positive term series $\sum_{n=1}^{\infty} a_n$ converges, then the series $\sum_{n=1}^{\infty} A_n$, as $A_n = \sqrt{\sum_{k=n}^{\infty} a_k} - \sqrt{\sum_{k=n+1}^{\infty} a_k}$, 
        also converges, and $a_n = o(A_n), n \to \infty$;
        \item[d)] If the positive term series $\sum_{n=1}^{\infty} a_n$ diverges, then the series $\sum_{n=1}^{\infty} A_n$, as $A_n = \sqrt{\sum_{k=1}^{n} a_k} - \sqrt{\sum_{k=1}^{n-1} a_k}$, also diverges, 
        and $A_n = o(a_n), n \to \infty$.
    \end{enumerate}
\end{prb}
\fbox{\begin{minipage}{0.95\textwidth}
\begin{sol*}
\begin{enumerate}
    \item[a)] Applying the Cauchy convergence criterion to the series $\sum_{n=1}^{\infty} a_n$, we have: $\forall \varepsilon > 0, \exists N > 0 \text{ such that } \forall n, m > N, |a_n + \cdots + a_m| < \varepsilon$. 
    If we set $m = 2n$ and since $a_n$ is monotonically decreasing, then $n a_{2n} < \varepsilon$, i.e. $a_n = o\left(\frac{1}{n}\right)$;
    \item[b)] Consider alternating series and monotonic sequences: $a_n \triangleq (-1)^n \sup_{k \geq n} b_k$, then $\sum_{n=1}^{\infty} a_n$ converges. (Leibniz Theorem) Or just let $a_n = (-1)^n \frac{1}{n}$.
    \item[c)] d$)$ Easy to prove. (Cancellation of intermediate terms in summation; rationalization, division).
\end{enumerate}
\end{sol*}
\end{minipage}}
\newline\newline
\fbox{\begin{minipage}{0.95\textwidth}
\textcolor{skycyan}{Thoughts:} \\
From c) and d), it can be seen: there always exists a convergent positive term series that is "larger" than a given convergent positive term series, and there always exists a positive term series that is "smaller" than 
a given divergent positive term series. \\ There is no unique, universal standard series that can be applied to all comparison tests.
\end{minipage}}

\begin{prb}
    (C3S2Q7\&Q8)
    \begin{enumerate}
        \item[a)] The series $\sum_{n=1}^{\infty} \ln a_n, a_n > 0, n \in \mathbb{N}$ converges i.f.f. the sequence $\{\Pi_n = a_1 \cdots a_n\}$ has a nonzero limit;
        \item[b)] If the infinite product $\prod_{n=1}^{\infty} e_n$ converges, then $\lim_{n \to \infty} e_n = 1$;
        \item[c)] The series $\sum_{n=1}^{\infty} \ln (1 + a_n), |a_n| < 1$ converges absolutely i.f.f. the series $\sum_{n=1}^{\infty} a_n$ converges absolutely;
        \item[d)] The infinite product $\prod_{n=1}^{\infty} e_n, e_n > 0$ converges i.f.f. the series $\sum_{n=1}^{\infty} \ln e_n$ converges;
        \item[e)] If $e_n = 1 + a_n$ and $a_n$ all have the same sign, then the infinite product $\prod_{n=1}^{\infty} (1 + a_n)$ converges i.f.f. the series $\sum_{n=1}^{\infty} a_n$ converges.
    \end{enumerate}
\end{prb}
\fbox{\begin{minipage}{0.95\textwidth}
\begin{qsol*}
\begin{enumerate}
    \item[a)] Easy to prove. $\sum_{n=1}^{\infty} \ln a_n = \ln \prod_{n=1}^{\infty} a_n$ (they have the same convergence behavior).
    \item[b)] Easy to prove. (If $\prod_{n=1}^{\infty} e_n = c \neq 0$, then $\lim_{n \to \infty} e_n = \lim_{n \to \infty} \frac{\prod_{k=1}^{n} e_k}{\prod_{k=1}^{n-1} e_k} = 1$. Alternatively, take logarithms and apply 
    the basic properties of series convergence.)
    \item[c)] If $\sum_{n=1}^{\infty} \ln(1 + a_n)$ converges, then by a) $\prod_{n=1}^{\infty}(1 + a_n)$ has a nonzero limit. Further, by b), $a_n \to 0$ as $n \to \infty$, and since $|\ln(1 + a_n)| \sim \ln(1 + |a_n|) 
    \sim |a_n|$, it follows (from C3S2Q5 a)) that the statement is proved.
    \item[d)] Easy to prove. (Same as a))
    \item[e)] $\ln(1 + a_n) \sim a_n$, so $\prod_{n=1}^{\infty}(1 + a_n)$ converges if and only if $\sum_{n=1}^{\infty} \ln(1 + a_n)$ converges if and only if $\sum_{n=1}^{\infty} a_n$ converges.
\end{enumerate}
\end{qsol*}
\end{minipage}}

\begin{prb}
    (C3S2Q9)
    \begin{enumerate}
        \item[a)] Calculate $\Pi_{n=1}^{\infty}(1+x^{2^{n-1}})$;
        \item[b)] Caculate $\Pi_{n=1}^{\infty} \cos\left( \frac{x}{2^{n}} \right)$, and prove $\frac{\pi}{2}=\frac{1}{\sqrt{ \frac{1}{2} }\cdot\sqrt{ \frac{1}{2}+\frac{1}{2}\cdot\sqrt{ \frac{1}{2} } }\cdot\sqrt{ 
        \frac{1}{2}+\frac{1}{2}\cdot\sqrt{ \frac{1}{2}+\frac{1}{2}\cdot \sqrt{ \frac{1}{2} } } }\cdots}$ (Vieta's formula);
        \item[c)] Find the $f(x)$, s.t. $f(0)=1,\ f(2x)=\cos^{2}x\cdot f(x),\ \lim_{ x \to 0 }f(x)=f(0)$.
    \end{enumerate}
\end{prb}
\fbox{\begin{minipage}{0.95\textwidth}
\begin{sol*}
\begin{enumerate}
    \item[a)] From C3S2Q7 e), $\prod_{n=1}^{\infty}(1+x^{2^n-1})$ and $\sum_{n=1}^{\infty}x^{2^n-1}$ have the same convergence behavior, i.e., they converge only when $|x|<1$, and $\prod_{n=1}^{\infty}(1+x^{2^n-1})=\frac{1-x^{2^n}}{1-x}$ 
    (generalized power series multiplication theorem);
    \item[b)] $\prod\limits_{n=1}^{\infty}\cos\left(\frac{x}{2^n}\right)=\lim\limits_{n\to\infty}\prod_{k=1}^{n}\cos\left(\frac{x}{2^k}\right)=\lim\limits_{n\to\infty}\frac{\sin\frac{x}{2^n}\prod_{k=1}^{n}\cos\frac{x}{2^k}}{\sin\frac{x}{2^k}}
    =\lim\limits_{n\to\infty}\frac{\sin x}{2^n\sin\frac{x}{2^n}}=\lim\limits_{n\to\infty}\frac{\sin x}{x}\cdot\frac{\frac{x}{2^n}}{\sin\frac{x}{2^n}}=\frac{\sin x}{x}$ Due to $\cos\frac{x}{2}=\sqrt{\frac{1+\cos x}{2}}$ and (above) 
    $\prod_{n=1}^{\infty}\cos\frac{x}{2^n}=\frac{\sin x}{x}$ (let $x=\frac{\pi}{2}$), Vieta's formula is proved;
    \item[c)] Using the recurrence relation $f(2x)=\cos^2 x \cdot f(x)$ and taking $\frac{x}{2^n}, n=0,1,2,\ldots$ yields, $f(x)=\left(\frac{\sin x}{x}\right)^2.$
\end{enumerate}
\end{sol*}
\end{minipage}}

\begin{prb}
    (C3S2Q10)
    \begin{enumerate}
        \item[a)] If $\frac{b_n}{b_{n+1}} = 1 + \beta_n, n = 1, 2, \cdots$, and the series $\sum_{n=1}^\infty \beta_n$ converges absolutely, then the limit $\lim_{n \to \infty} b_n = b$ evidently exists;
        \item[b)] If $\frac{a_n}{a_{n+1}} = 1 + \frac{p}{n} + \alpha_n, n = 1, 2, \cdots$, and the series $\sum_{n=1}^\infty \alpha_n$ converges absolutely, then $a_n \sim \frac{c}{n^p}, n \to \infty$;
        \item[c)] If $\frac{a_n}{a_{n+1}} = 1 + \frac{p}{n} + \alpha_n, n = 1, 2, \cdots$, and the series $\sum_{n=1}^\infty \alpha_n$ converges absolutely, then the series $\sum_{n=1}^\infty a_n$ converges when $p > 1$ and diverges 
        when $p \leq 1$ (Gauss's test for absolute convergence of series).
    \end{enumerate}
\end{prb}
\fbox{\begin{minipage}{0.95\textwidth}
\begin{sol*}
\begin{enumerate}
    \item[a)] From C3S2Q7 and Q8 c), it follows that $\sum_{n=1}^{\infty} \ln(1 + \beta_n) = \sum_{n=1}^{\infty} \ln \left( \frac{b_n}{b_{n+1}} \right)$ and $\sum_{n=1}^{\infty} \beta_n$ have the same convergence behavior, both absolutely convergent.
    \item[b)] (Good question, good solution) \\ Method 1. Consider $\frac{n^p}{(n+1)^p} \frac{a_n}{a_{n+1}} = \frac{n^p}{(n+1)^p} \left( 1 + \frac{p}{n} + \alpha_n \right) = \left( 1 + \frac{1}{n} \right)^{-p} + \left( 1 + \frac{1}{n} 
    \right)^{-p} \frac{p}{n} + \left( 1 + \frac{1}{n} \alpha_n \right)^{-p} = \left( 1 - \frac{p}{n} + O\left( \frac{1}{n^2} \right) \right) \left( 1 + \frac{p}{n} \right) + \left( 1 - \frac{p}{n} + O\left( \frac{1}{n^2} \right) \right) \alpha_n 
    \left[ \triangleq 1 + \theta_n \right]$, \\ $\theta_n = \left( 1 - \frac{p}{n} + O\left( \frac{1}{n^2} \right) \right) \left( 1 + \frac{p}{n} + \alpha_n \right) - 1 = O \left( \frac{1}{n^2} \right) \left( 1 + \frac{p}{n} + \alpha_n \right) - 
    \frac{p^2}{n^2} + \alpha_n \left( 1 - \frac{p}{n} \right)$, $|\theta_n| \leq |O \left( \frac{1}{n^2} \right) \left( 1 + \frac{p}{n} + \alpha_n \right)| + |\frac{p^2}{n^2}| + |\alpha_n \left( 1 - \frac{p}{n} \right)|$, and $\sum_{n=1}^{\infty} 
    \alpha_n$ abs. cgt. , then $\sum_{n=1}^{\infty} \theta_n$ abs. cgt. Then due to C3S2Q7 c), $\sum_{n=1}^{\infty} \ln(1 + \theta_n) = \sum_{n=1}^{\infty} \ln \frac{n^p a_n}{(n+1)^p a_{n+1}} = \ln \alpha_1 - 
    \lim_{n \to \infty} n^p \alpha_n$ abs. cgt. i.e. $n^p \alpha_n \sim c, n \to \infty.$ $(\alpha_n \sim \frac{c}{n^p}).$ \\ 
    Method 2 (Optimize). Following the first step of Method 1: \\ 
    $\ln \frac{n^p a_n}{(n+1)^p a_{n+1}} = \ln \frac{\alpha_n}{\alpha_{n+1}} - p \ln \left( 1 + \frac{1}{n} \right)$ $= \ln \left( 1 + \frac{p}{n} + \alpha_n \right) - p \ln \left( 1 + \frac{1}{n} \right)$ $= \frac{p}{n} + \alpha_n + O \left( 
    \frac{p}{n} + \alpha_n \right)^2 - p \left( \frac{1}{n} + O \left( \frac{1}{n^2} \right) \right)$ $= \alpha_n + O \left( \frac{1}{n^2} \right) + O (\alpha_n) \to$ abs. cgt. going on Method 1.
    \item[c)] From b) and the convergence behavior of $\sum_{n=1}^{\infty} \frac{1}{n^p}$, it immediately follows.
\end{enumerate}
\end{sol*}
\end{minipage}}

\begin{prb}
    (C3S2Q11) \\
    Prove: For any sequence $\{a_n\}$, $\limsup \left( \frac{a_{n+1}}{a_n} \right)^n > e$, and this estimate is the best possible.
\end{prb}
\fbox{\begin{minipage}{0.95\textwidth}
\begin{sol*}
(Great deal: Wonderful observation!) /Contradiction/ \\
Lemma: $\forall x\in(0,e)$, for sufficiently large $n$, $x^{\frac{1}{n}}<1+\frac{1}{n}$. \\
If $\varlimsup\limits_{ n \to \infty }\left( \frac{1+a_{n+1}}{a_{n}} \right)^{n}<e,\ a_{n}>0$, then $\exists b$, s.t. $\varlimsup\limits_{ n \to \infty }\left( \frac{1+a_{n+1}}{a_{n}} \right)^{n}<b<e$,
i.e. $\exists$ sufficient large $n$, s.t. $\left( \frac{1+a_{n+1}}{a_{n}} \right)^{n}<b<e$,
i.e. (using std::Lemma) $\frac{1+a_{n+1}}{a_{n}}<b^{\frac{1}{n}}<\frac{n+1}{n}$,
let $b_{n} \triangleq \frac{a_{n}}{n}$, then $b_{n+1}<b_{n}-\frac{1}{n+1}$, increasing to $n+m$ and $\Sigma$ them: $b_{n+m}<b_{n}-\sum_{k=1}^{m} \frac{1}{n+k}$, fixing $b_{n}$ and letting $m\to \infty$ we have $b_{n+m}=\frac{a_{n+m}}{n}\to {-}\infty,
\ m\to \infty$ which derivate constraction.
\end{sol*}
\end{minipage}}

%\end{document}